\documentclass[11pt]{article}

\usepackage[margin=1in]{geometry}
\usepackage{amsmath,amssymb,amsthm,mathtools}
\usepackage{booktabs,array}
\usepackage{microtype}
\setlength{\emergencystretch}{1em}
\usepackage{enumitem}
\usepackage[hidelinks]{hyperref}

\hypersetup{
  pdftitle={Quadratic Line-Cover Certificates and the 2n-4 Bound for Checkerboard No-Three-in-Line Sets},
  pdfauthor={Dominic A. Dabish},
  pdfsubject={Checkerboard no-three-in-line sets; quadratic line-cover certificates; Lean 4 formalization}
}

\newtheorem{theorem}{Theorem}[section]
\newtheorem{lemma}[theorem]{Lemma}
\newtheorem{proposition}[theorem]{Proposition}
\newtheorem{corollary}[theorem]{Corollary}
\theoremstyle{remark}
\newtheorem{remark}[theorem]{Remark}

\newcommand{\Dmono}{D_{\mathrm{mono}}}
\newcommand{\Lmono}{L_{\mathrm{mono}}}

\title{Quadratic Line-Cover Certificates and the $2n-4$ Bound\\for Checkerboard No-Three-in-Line Sets}
\author{Dominic A. Dabish\\\small San Diego State University}
\date{July 2026}

\begin{document}
\maketitle

\begin{abstract}
Let $\Dmono(n)$ be the largest number of points that can be chosen from one parity class of the $n\times n$ integer grid with no three collinear. Prellberg proved the elementary bound $\Dmono(n)\le 2n-2$ and recorded the sharper inequality $\Dmono(n)\le 2n-4$ for $n\ge6$ as a conjectural consequence of a boundary-forcing mechanism. We prove this inequality for every $n\ge6$.

For all odd and even boards outside two finite exceptions, the proof gives explicit nonnegative quadratic weights on the rows, columns, and the two diagonal families of slopes $\pm1$. Every point in the prescribed parity class receives the same total weight, while the exact capacity-two cost is strictly below the threshold required by $2n-3$ selected points. The thin colour of the $7\times7$ board is handled by a small integer line cover. The two $6\times6$ colour classes are handled by an exact finite Boolean verification. A complete Lean~4 formalization proves the all-$n$ theorem in the genuine all-slope no-three-in-line model and includes an explicit axiom audit.
\end{abstract}

\noindent\textbf{Mathematics Subject Classifications (2020):} 52C10, 05B40, 90C05, 03B35.\\
\textbf{Keywords:} no-three-in-line problem, checkerboard grid, line-cover certificate, linear-programming duality, formal proof, Lean.

\section{Introduction}

The classical no-three-in-line problem asks for the largest number of points that can be chosen from an $n\times n$ integer grid with no three collinear. The problem goes back at least to Dudeney's 1917 chessboard puzzle~\cite{dudeney}. Its systematic study includes the foundational work of Guy and Kelly~\cite{guy-kelly}, Adena, Holton, and Kelly~\cite{adena-holton-kelly}, and Hall, Jackson, Sudbery, and Wild~\cite{hall-jackson-sudbery-wild}.

The ordinary problem has developed along both structural and computational lines. Flammenkamp's searches produced extensive families of $2n$-point configurations and complete enumerations for small boards~\cite{flammenkamp-1992,flammenkamp-1998}; more recently, constraint-satisfaction methods extended the known equality $D(n)=2n$ through $n=60$~\cite{prellberg-csp}. Related directions include dense extensible subsets of the infinite grid~\cite{nagy-nagy-woodroofe} and the no-$(k+1)$-in-line problem, which has now been resolved for every fixed $k\ge3$ and all sufficiently large $n$~\cite{ghosal-et-al}. These developments emphasize the exceptional difficulty of the original capacity-two case.

Write
\[
G_n=\{0,1,\ldots,n-1\}^2
\]
and, for $\varepsilon\in\{0,1\}$, let
\[
C_\varepsilon(n)=\{(x,y)\in G_n:x+y\equiv\varepsilon\pmod 2\}.
\]
A set is \emph{no-three-in-line} (NTIL) if it contains no three collinear points. Define
\[
\Dmono(n,\varepsilon)
 =\max\{|S|:S\subseteq C_\varepsilon(n)\text{ is NTIL}\},
\qquad
\Dmono(n)=\max_{\varepsilon\in\{0,1\}}\Dmono(n,\varepsilon).
\]

Prellberg~\cite{prellberg-checkerboard} introduced this fixed-parity variant and observed that the monochromatic diagonals of either slope $+1$ or slope $-1$ give
\[
\Dmono(n)\le 2n-2.
\]
He also recorded a boundary-forcing mechanism that suggests, but does not prove,
\[
\Dmono(n)\le 2n-4\qquad(n\ge6).
\]
His main framework is a four-direction linear-programming relaxation, together with symmetry reductions and a continuum dual problem. We prove the stated finite upper bound by constructing closed-form discrete dual certificates and treating the two residual board sizes exactly.

The proof uses only a necessary consequence of the full NTIL condition: every row, column, slope-$1$ diagonal, and slope-$(-1)$ diagonal contains at most two selected points. A certificate using these four line families therefore gives an upper bound for the original all-slope problem; it does not replace collinearity by a four-direction definition.

\begin{theorem}[Main theorem]\label{thm:main}
For every integer $n\ge6$ and each $\varepsilon\in\{0,1\}$,
\[
\Dmono(n,\varepsilon)\le 2n-4.
\]
Consequently $\Dmono(n)\le2n-4$ for every $n\ge6$.
\end{theorem}

The result is complementary to the continuum analysis in~\cite{prellberg-checkerboard}. Rather than passing to a scaling limit, we exhibit finite-board weights whose coverage and cost can be evaluated exactly. For $n\ge7$ these weights certify a strict fractional inequality; only the $6\times6$ board requires an additional integrality argument.

The proof is elementary once the quadratic weights are found. The identity behind them is
\begin{equation}\label{eq:master-identity}
(2x-N)^2+(2y-N)^2
 +2\bigl(N^2-(x+y-N)^2\bigr)
 +2\bigl(N^2-(x-y)^2\bigr)=4N^2.
\end{equation}
The parity restriction selects exactly the two diagonal weights that occur at a point. Exact sums of squares then give the cost.

\section{A weighted line-cover lemma}

Let $\mathcal L$ consist of all rows, columns, and diagonals of slopes $\pm1$ that meet $G_n$. Assign a nonnegative weight $w_\ell$ to every $\ell\in\mathcal L$. For $p\in G_n$, define
\[
\operatorname{cov}(p)=\sum_{\ell\ni p}w_\ell,
\qquad
\operatorname{cost}(w)=2\sum_{\ell\in\mathcal L}w_\ell.
\]
The factor $2$ is the NTIL capacity of each line.

\begin{lemma}[Line-cover certificate]\label{lem:certificate}
Suppose $S\subseteq G_n$ is NTIL and every $p\in S$ satisfies $\operatorname{cov}(p)\ge q$, where $q>0$. Then
\[
q|S|\le \operatorname{cost}(w).
\]
In particular, if $\operatorname{cost}(w)<q(k+1)$, then $|S|\le k$.
\end{lemma}

\begin{proof}
Double-count weighted incidences:
\[
q|S|
 \le \sum_{p\in S}\operatorname{cov}(p)
 =\sum_{\ell\in\mathcal L}w_\ell|S\cap\ell|
 \le2\sum_{\ell\in\mathcal L}w_\ell.
\]
The last assertion follows because $|S|$ is an integer.
\end{proof}

Let $\Lmono(n,\varepsilon)$ denote the optimum of Prellberg's four-direction linear-programming relaxation. After division by $q$, the weights in Lemma~\ref{lem:certificate} are precisely a feasible solution of its dual~\cite{prellberg-checkerboard,schrijver}. For $n\ge7$, our certificates prove the stronger fractional statement
\[
\Lmono(n,\varepsilon)<2n-3,
\]
and the desired integral bound follows by rounding down. The $6\times6$ board is the sole case in which an additional integrality argument is required.

\section{Odd side length}

Let $n=2m+1$ and put $N=2m$. The class containing $(0,0)$ is the \emph{fat} class, and the other is the \emph{thin} class.

For every row and column with index $i\in\{0,\ldots,2m\}$, assign weight
\[
R_i=(2i-2m)^2.
\]
For a diagonal indexed by $j\in\{0,\ldots,4m\}$, assign
\[
B_j=2\bigl((2m)^2-(j-2m)^2\bigr)
\]
when $j$ has the parity of the chosen colour, and weight $0$ otherwise. The two diagonal indices through $(x,y)$ are $j=x+y$ and $j=x-y+2m$. Since $2m$ is even, both have parity $x+y$.

All weights are nonnegative. For every point $(x,y)$ in the chosen parity class, the four incident weights sum, by~\eqref{eq:master-identity}, to
\[
q=16m^2.
\]

The needed exact sums are
\begin{align}
A_m&:=\sum_{i=0}^{2m}(2i-2m)^2
 =\frac{4m(m+1)(2m+1)}3,\label{eq:odd-axis}\\
F_m&:=\sum_{\substack{0\le j\le4m\\j\equiv0\pmod2}}B_j
 =\frac{8m(2m-1)(2m+1)}3,\label{eq:odd-fat-diag}\\
T_m&:=\sum_{\substack{0\le j\le4m\\j\equiv1\pmod2}}B_j
 =\frac{4m(8m^2+1)}3.\label{eq:odd-thin-diag}
\end{align}
Derivations are included in Appendix~\ref{app:sums}.

\subsection{The fat class}

There are two axis families and two diagonal families, so the capacity-two cost is
\[
C_{\mathrm{fat}}=2(2A_m+2F_m)
 =\frac{16m(10m^2+3m-1)}3.
\]
The target is $|S|\le4m-2$. We compute
\[
q(4m-1)-C_{\mathrm{fat}}
 =\frac{16m}{3}(2m^2-6m+1).
\]
For $m\ge3$,
\[
2m^2-6m+1=2m(m-3)+1>0.
\]
Lemma~\ref{lem:certificate} therefore gives
\[
|S|\le4m-2=2n-4
\]
for every odd fat board with $n\ge7$.

\subsection{The thin class}

The cost is
\[
C_{\mathrm{thin}}=2(2A_m+2T_m)
 =\frac{16m(10m^2+3m+2)}3.
\]
For $m\ge4$,
\[
q(4m-1)-C_{\mathrm{thin}}
 =\frac{16m}{3}(2m^2-6m-2)>0,
\]
because the quadratic factor is increasing for $m\ge4$ and equals $6$ at $m=4$. Hence
\[
|S|\le4m-2=2n-4
\]
for every odd thin board with $n\ge9$.

\subsection{The exceptional thin \texorpdfstring{$7\times7$}{7 x 7} board}

It remains to treat $m=3$ in the thin class. Give weight $1$ to each boundary row or column with index $0$ or $6$. For each of the two diagonal families, indexed by $j\in\{0,\ldots,12\}$, define
\[
f(j)=
\begin{cases}
2,&|j-6|=1,\\
1,&|j-6|=3,\\
0,&\text{otherwise}.
\end{cases}
\]
Only odd $j$ occur in the thin class. The total line-weight sum is $16$, hence the capacity-two cost is $32$. Up to the dihedral symmetries of the square, the $24$ thin-colour points form four orbits:

\begin{center}
\begin{tabular}{ccc}
\toprule
Representative & Orbit size & Coverage\\
\midrule
$(0,1)$ & $8$ & $3$\\
$(0,3)$ & $4$ & $3$\\
$(1,2)$ & $8$ & $3$\\
$(2,3)$ & $4$ & $4$\\
\bottomrule
\end{tabular}
\end{center}
Thus every thin point has coverage at least $q=3$, and
\[
32<3\cdot11.
\]
Lemma~\ref{lem:certificate} gives $|S|\le10=2\cdot7-4$.

\section{Even side length}

Let $n=2m$ and put $N=2m-1$. Assign to each row and column index $i\in\{0,\ldots,2m-1\}$ the weight
\[
R_i=(2i-N)^2,
\]
and define, for $j\in\{0,\ldots,2N\}$,
\[
B_j=2\bigl(N^2-(j-N)^2\bigr).
\]
On the sum diagonals $j=x+y$, retain the parity of the chosen colour. On the shifted difference diagonals $j=x-y+N$, retain the opposite parity, because $N$ is odd.

Every point of the chosen colour receives exactly the four terms in~\eqref{eq:master-identity}, so its coverage is
\[
q=4N^2.
\]
The axis sum and total selected diagonal sum are
\begin{align*}
A'_m&=\sum_{i=0}^{2m-1}(2i-N)^2
 =\frac{2m(4m^2-1)}3,\\
D'_m&=\sum_{j=0}^{2N}B_j
 =\frac{2N(4m-3)(4m-1)}3.
\end{align*}
The two chosen diagonal parity classes, one in each diagonal family, partition the indices $0,\ldots,2N$. Hence
\[
C_{\mathrm{even}}=2(2A'_m+D'_m)
 =\frac{4N(20m^2-14m+3)}3.
\]
For $m\ge4$,
\begin{align*}
q(4m-3)-C_{\mathrm{even}}
 &=\frac{8(2m-1)}3(2m^2-8m+3)>0,
\end{align*}
because $2m^2-8m+3=2m(m-4)+3$. Therefore
\[
|S|\le4m-4=2n-4
\]
for every even board with $n\ge8$.

\section{The \texorpdfstring{$6\times6$}{6 x 6} base case}\label{sec:n6}

Each parity class of the $6\times6$ board contains $18$ points. At $m=3$, the even quadratic certificate above has normalized cost
\[
\frac{C_{\mathrm{even}}}{q}=\frac{940}{100}=\frac{47}{5},
\]
so it does not exclude nine selected points. A finite integrality argument completes the proof.

\begin{proposition}\label{prop:n6}
Every subset of either parity class of the $6\times6$ board that meets every row, column, and slope-$\pm1$ diagonal in at most two points has size at most $8$.
\end{proposition}

\begin{proof}[Exact finite verification]
Fix a parity class and enumerate its $18$ points. It is enough to reject all nine-point subsets: any larger feasible set would contain a feasible nine-point subset. There are
\[
\binom{18}{9}=48{,}620
\]
such subsets. For each subset, the checker evaluates every row, column, sum diagonal, and difference diagonal and rejects the subset as soon as one line has at least three selected points. No nine-point subset survives. As a reproducibility checksum, exactly $155$ eight-point subsets survive in each parity class under these four line-capacity constraints.

The accompanying standard-library Python program performs this enumeration directly. Independently, the Lean development transports an arbitrary monochromatic set to $18$ Boolean variables and proves the same statement with the kernel-checked bit-vector procedure \texttt{bv\_decide}. Since every all-slope NTIL set satisfies the four line-capacity constraints, the proposition follows.
\end{proof}

\begin{proof}[Proof of Theorem~\ref{thm:main}]
The odd fat case holds for $n\ge7$, the odd thin case for $n\ge9$ with the separate certificate at $n=7$, and the even case for $n\ge8$. Proposition~\ref{prop:n6} handles $n=6$.
\end{proof}

\section{Formal verification and reproducibility}

The complete theorem is formalized in Lean~4 with Mathlib~\cite{lean4,mathlib}. This use of a small trusted proof kernel follows the broader tradition of machine-checked discrete mathematics exemplified by the formal proofs of the Four Color Theorem and the Kepler conjecture~\cite{gonthier-four-color,hales-flyspeck}. The public declaration is
\begin{verbatim}
theorem Checkerboard.checkerboard_upper_all_n {n parity : Nat}
    (hn : 6 <= n) (hp : Or (parity = 0) (parity = 1))
    (s : Finset (Point n))
    (hcolor : Monochromatic parity s)
    (hntil : NoThreeInLine s) :
    s.card <= 2 * n - 4
\end{verbatim}
The predicate \texttt{NoThreeInLine} quantifies over every triple of distinct selected points and uses the exact integer determinant. Thus the formal theorem is the genuine all-slope statement. The four principal line capacities are proved from that hypothesis.

The principal formal declarations corresponding to the paper are:
\begin{center}
\begin{tabular}{>{\raggedright\arraybackslash}p{0.42\textwidth}>{\raggedright\arraybackslash}p{0.48\textwidth}}
\toprule
Mathematical component & Lean declaration\\
\midrule
Weighted double-counting certificate & \texttt{card\_le\_of\_fourCertificate}\\
Odd/even quadratic weights & \texttt{oddQuadraticWeights}, \texttt{evenQuadraticWeights}\\
Constant point coverage & \texttt{oddQuadratic\_coverage}, \texttt{evenQuadratic\_coverage}\\
Exact certificate costs & \texttt{oddQuadratic\_cost\_zero}, \texttt{oddQuadratic\_cost\_one}, \texttt{evenQuadratic\_cost\_zero}, \texttt{evenQuadratic\_cost\_one}\\
Odd fat and thin bounds & \texttt{odd\_zero\_upper}, \texttt{odd\_one\_upper}\\
Even bound & \texttt{even\_upper}\\
Thin $7\times7$ exception & \texttt{n7\_one\_upper}\\
Two $6\times6$ base cases & \texttt{n6\_zero\_upper}, \texttt{n6\_one\_upper}\\
All boards $n\ge6$ & \texttt{checkerboard\_upper\_all\_n}\\
\bottomrule
\end{tabular}
\end{center}

The project is pinned to Lean~4.32.0. Its audit workflow rejects unfinished proof markers, project-defined axioms, \texttt{native\_decide}, opaque shortcuts, and any occurrence of \texttt{sorryAx} in the compiled output. The axiom audit explicitly runs
\begin{verbatim}
#print axioms Checkerboard.checkerboard_upper_all_n
\end{verbatim}
The repository records the exact commit, toolchain, workflow run, and independent finite checkers.

\section{Scope and open questions}

Theorem~\ref{thm:main} proves the finite upper bound proposed in Prellberg's paper. It does not determine $\Dmono(n)$ exactly for every $n$, and it does not settle the conjectured asymptotic value
\[
\lim_{n\to\infty}\frac{\Dmono(n)}n.
\]
The present line covers use only four necessary line families. Proving a matching all-slope lower bound requires controlling every rational slope and is a separate problem.

To the author's knowledge, this is the first proof of the bound $\Dmono(n)\le2n-4$ for all $n\ge6$. This priority statement is intentionally qualified because unindexed or simultaneous work may exist.

\section*{Data and code availability}

The Lean formalization, independent Python checkers, manuscript source, and reproducibility record are contained in the public repository \url{https://github.com/DomTheDeveloper/crl}. The proof branch and exact commit identifier are recorded in the accompanying reproducibility file. An archival release should be deposited before journal submission.

\appendix
\section{Derivation of the line-weight sums}\label{app:sums}

We record the elementary calculations used above. First,
\[
\sum_{r=1}^{m}r^2=\frac{m(m+1)(2m+1)}6,
\qquad
\sum_{r=1}^{m}(2r-1)^2=\frac{m(4m^2-1)}3.
\]
For the odd axis weights,
\[
A_m=4\sum_{r=-m}^{m}r^2
 =8\sum_{r=1}^{m}r^2
 =\frac{4m(m+1)(2m+1)}3.
\]
For even diagonal indices $j=2r$,
\begin{align*}
F_m
 &=8\sum_{r=0}^{2m}\bigl(m^2-(r-m)^2\bigr)\\
 &=8\left((2m+1)m^2-\frac{m(m+1)(2m+1)}3\right)\\
 &=\frac{8m(2m-1)(2m+1)}3.
\end{align*}
For odd diagonal indices, the centered values $j-2m$ are the odd integers from $-(2m-1)$ to $2m-1$. Therefore
\begin{align*}
T_m
 &=2\left(2m\cdot4m^2-2\sum_{r=1}^{m}(2r-1)^2\right)\\
 &=\frac{4m(8m^2+1)}3.
\end{align*}

For the even board, the axis coordinates are the odd integers from $-(2m-1)$ to $2m-1$, giving
\[
A'_m=2\sum_{r=1}^{m}(2r-1)^2
 =\frac{2m(4m^2-1)}3.
\]
Finally, with $N=2m-1$,
\begin{align*}
D'_m
 &=2\sum_{t=-N}^{N}(N^2-t^2)\\
 &=2\left((2N+1)N^2-\frac{N(N+1)(2N+1)}3\right)\\
 &=\frac{2N(2N-1)(2N+1)}3
 =\frac{2N(4m-3)(4m-1)}3.
\end{align*}

\begingroup
\raggedright
\begin{thebibliography}{99}

\bibitem{dudeney}
H. E. Dudeney,
\newblock \emph{Amusements in Mathematics},
\newblock Thomas Nelson and Sons, London, 1917, Puzzle 317.

\bibitem{guy-kelly}
R. K. Guy and P. A. Kelly,
\newblock The no-three-in-line problem,
\newblock \emph{Canadian Mathematical Bulletin} 11 (1968), 527--531,
\newblock \href{https://doi.org/10.4153/CMB-1968-062-3}{doi:10.4153/CMB-1968-062-3}.

\bibitem{adena-holton-kelly}
M. A. Adena, D. A. Holton, and P. A. Kelly,
\newblock Some thoughts on the no-three-in-line problem,
\newblock in \emph{Combinatorial Mathematics}, Lecture Notes in Mathematics 403,
Springer, 1974, pp. 6--17.

\bibitem{hall-jackson-sudbery-wild}
R. R. Hall, T. H. Jackson, A. Sudbery, and K. Wild,
\newblock Some advances in the no-three-in-line problem,
\newblock \emph{Journal of Combinatorial Theory, Series A} 18 (1975), 336--341,
\newblock \href{https://doi.org/10.1016/0097-3165(75)90043-6}{doi:10.1016/0097-3165(75)90043-6}.

\bibitem{flammenkamp-1992}
A. Flammenkamp,
\newblock Progress in the no-three-in-line problem,
\newblock \emph{Journal of Combinatorial Theory, Series A} 60 (1992), 305--311,
\newblock \href{https://doi.org/10.1016/0097-3165(92)90012-J}{doi:10.1016/0097-3165(92)90012-J}.

\bibitem{flammenkamp-1998}
A. Flammenkamp,
\newblock Progress in the no-three-in-line problem, II,
\newblock \emph{Journal of Combinatorial Theory, Series A} 81 (1998), 108--113,
\newblock \href{https://doi.org/10.1006/jcta.1997.2829}{doi:10.1006/jcta.1997.2829}.

\bibitem{prellberg-csp}
T. Prellberg,
\newblock Constraint satisfaction programming for the no-three-in-line problem,
\newblock arXiv:2602.07751, 2026,
\newblock \href{https://doi.org/10.48550/arXiv.2602.07751}{doi:10.48550/arXiv.2602.07751}.

\bibitem{nagy-nagy-woodroofe}
D. T. Nagy, Z. L. Nagy, and R. Woodroofe,
\newblock The extensible no-three-in-line problem,
\newblock \emph{European Journal of Combinatorics} 114 (2023), Paper 103796,
\newblock \href{https://doi.org/10.1016/j.ejc.2023.103796}{doi:10.1016/j.ejc.2023.103796}.

\bibitem{ghosal-et-al}
A. Ghosal, R. Goenka, A. Grebennikov, P. Keevash, M. Kwan, and H. T. Pham,
\newblock No-$(k+1)$-in-line problem for $k\ge3$,
\newblock arXiv:2607.05255, 2026,
\newblock \href{https://doi.org/10.48550/arXiv.2607.05255}{doi:10.48550/arXiv.2607.05255}.

\bibitem{prellberg-checkerboard}
T. Prellberg,
\newblock No-three-in-line sets on the checkerboard grid,
\newblock arXiv:2605.09215, 2026,
\newblock \href{https://doi.org/10.48550/arXiv.2605.09215}{doi:10.48550/arXiv.2605.09215}.

\bibitem{schrijver}
A. Schrijver,
\newblock \emph{Theory of Linear and Integer Programming},
\newblock Wiley, 1986.

\bibitem{lean4}
L. de Moura and S. Ullrich,
\newblock The Lean 4 theorem prover and programming language,
\newblock in \emph{Automated Deduction -- CADE 28}, Lecture Notes in Computer Science 12699,
Springer, 2021, pp. 625--635,
\newblock \href{https://doi.org/10.1007/978-3-030-79876-5_37}{doi:10.1007/978-3-030-79876-5\_37}.

\bibitem{mathlib}
The Mathlib Community,
\newblock The Lean mathematical library,
\newblock in \emph{Proceedings of the 9th ACM SIGPLAN International Conference on Certified Programs and Proofs}, 2020, pp. 367--381,
\newblock \href{https://doi.org/10.1145/3372885.3373824}{doi:10.1145/3372885.3373824}.

\bibitem{gonthier-four-color}
G. Gonthier,
\newblock Formal proof---the Four Color Theorem,
\newblock \emph{Notices of the American Mathematical Society} 55 (2008), no. 11, 1382--1393.

\bibitem{hales-flyspeck}
T. Hales, M. Adams, G. Bauer, T. D. Dang, J. Harrison, L. T. Hoang,
C. Kaliszyk, V. Magron, S. McLaughlin, T. T. Nguyen, Q. T. Nguyen,
T. Nipkow, S. Obua, J. Pleso, J. Rute, A. Solovyev, T. H. A. Ta,
N. T. Tran, T. D. Trieu, J. Urban, K. Vu, and R. Zumkeller,
\newblock A formal proof of the Kepler conjecture,
\newblock \emph{Forum of Mathematics, Pi} 5 (2017), e2,
\newblock \href{https://doi.org/10.1017/fmp.2017.1}{doi:10.1017/fmp.2017.1}.

\end{thebibliography}
\endgroup

\section*{Declaration of generative AI and AI-assisted technologies in the manuscript preparation process}

During the preparation of this work, the author used OpenAI ChatGPT to assist with proof exploration, development and review of Lean and Python code, literature organization, and improvements to the clarity and structure of the manuscript. After using this tool, the author reviewed and edited all material incorporated into the work, checked the cited sources, and takes full responsibility for the mathematical content and the final text. The system used was OpenAI ChatGPT with the GPT-5.6 Thinking model (accessed July 2026).

\end{document}
